%!TEX root = ../sztuthesis_main.tex
% 设置中文摘要

\phantomsection
\addcontentsline{toc}{section}{摘要}

\keywordscn{大语言模型；个人知识库；智能优化；知识图谱}
\categorycn{TP391}
\begin{abstractcn}
  在个人博客，技术笔记等数字化个人文档资料管理的场景下，用户会积累越来越多的文档资料。但这些高价值的文本信息一直缺乏一种能够适用于长期积累与较大规模数据场景下的有效管理和汇聚方法。随着记录的博客和笔记数量越来越多，这些文档资料往往会变得越来越零散，甚至会被用户遗忘或丢弃掉，因为对于用户来说，随着博客和笔记数量的不断增加，人工手动整理和维护的成本也不断增加。

  为了解决大量的个人非结构化文本资料的管理难题，本文设计并实现了一个由大语言模型驱动的个人知识库构建与优化系统 MindTrace。能够把原始文档信息，即用户的博客，笔记或任何其它文本资料，通过大语言模型处理流水线把用户原始文档资料沉淀为一个个实体（Entity），观念（Concept），以及主题（Topic）。并且依托异步任务解耦技术，让用户无感于后台基于 LLM 的文档数据处理流水线。

  结果表明，MindTrace 在个人文档管理场景下是行之有效的，系统根据文档中的观念和主题自行智能归类，免去人工自行管理的烦恼，并从看似零散的个人文档资料中提取出有价值的信息，进一步识别文档与文档之间的观念，主题联系，得到用户长期形成的观念脉络，为用户长期知识沉淀提供一种解决办法。
\end{abstractcn}
